\section{Quotient-cover obstruction and refined constants}
\label{sec:quotient-refinements}

The fixed-label correction has consequences which are not visible from a
single character channel.  We first identify the exact geometric factor
missed by the Artin substitution and then refine the product asymptotic.

\subsection{Normal quotients and the missing split-orbit factor}

Let \(H\triangleleft G\), put \(Q=G/H\) and \(q=|Q|\), and write
\(\tau_Q:E\to Q\) for the quotient cocycle.  The adjacency matrix of the
regular \(Q\)-skew product is denoted by \(\widetilde A_Q\); explicitly, its
vertices are \(V\times Q\), and an edge \(e:i\to j\) gives an edge
\((i,a)\to(j,a\tau_Q(e))\) for each \(a\in Q\).  Define
\begin{equation}\label{eq:quotient-cover-reduced-constant}
 C(\widetilde A_Q)
 :=\lim_{t\uparrow1}(1-t)
       \det(I-t\widetilde A_Q/\lambda)^{-1}.
\end{equation}
Under the strict twisted gap this limit is positive: the regular
representation decomposition shows that the only peripheral block is the
simple Perron block of \(A\).

For a primitive orbit \(\omega\), let
\begin{equation}\label{eq:quotient-holonomy-order}
 q_H(\omega):=\operatorname{ord}_Q(g_\omega H).
\end{equation}
This is independent of the based representative.  Finally, put
\begin{equation}\label{eq:normal-subgroup-product-constant}
 K_H:=\prod_{\substack{C\in\operatorname{Conj}(G)\\C\subset H}}K_C.
\end{equation}

\begin{theorem}[Exact quotient-cover splitting factor]
\label{thm:quotient-cover-splitting}
Assume the hypotheses of \Cref{thm:correct-product}.  For every normal
subgroup \(H\triangleleft G\),
\begin{equation}\label{eq:quotient-cover-splitting}
 K_H=
 \left(\frac{e^{-\gamma}}{C(\widetilde A_Q)}\right)^{1/q}
 \prod_{\substack{\omega\in\cP\\q_H(\omega)>1}}
 \left(1-\lambda^{-q_H(\omega)|\omega|}\right)^{-1/q_H(\omega)}.
\end{equation}
The product in \eqref{eq:quotient-cover-splitting} converges absolutely.
If \(K_H^{\mathrm{Art}}\) denotes the candidate obtained by replacing every
fixed-label coordinate by its periodic Artin coordinate, then
\begin{equation}\label{eq:quotient-artin-candidate}
 K_H^{\mathrm{Art}}
 =\left(\frac{e^{-\gamma}}{C(\widetilde A_Q)}\right)^{1/q}
\end{equation}
and
\begin{equation}\label{eq:quotient-artin-positive-error}
 \log\frac{K_H}{K_H^{\mathrm{Art}}}
 =-\sum_{\substack{\omega\in\cP\\q_H(\omega)>1}}
    \frac{1}{q_H(\omega)}
    \log\left(1-\lambda^{-q_H(\omega)|\omega|}\right).
\end{equation}
Consequently, the Artin candidate is strictly smaller than \(K_H\) if and
only if some primitive orbit has \(g_\omega\notin H\).
\end{theorem}

\begin{proof}
Inflate each irreducible representation of \(Q\) to \(G\).  The regular
character identity gives
\begin{equation}\label{eq:normal-subgroup-indicator}
 \mathbf1_H(g)=\frac{1}{q}\sum_{\sigma\in\Irr(Q)}
       (\dim\sigma)\chi_\sigma(gH).
\end{equation}
Summing \eqref{eq:correct-product-fourier} over the conjugacy classes
contained in \(H\), or equivalently applying
\eqref{eq:normal-subgroup-indicator}, yields
\begin{equation}\label{eq:normal-subgroup-fixed-coordinate}
 \log K_H=-\frac{1}{q}\left(
   \gamma+\log C(A)+
   \sum_{\substack{\sigma\in\Irr(Q)\\\sigma\ne\mathbf1}}
      (\dim\sigma)F_{\chi_\sigma}(\lambda^{-1})\right).
\end{equation}

The regular representation decomposition also gives the polynomial identity
\begin{equation}\label{eq:quotient-regular-determinant-factorization}
 \det(I-z\widetilde A_Q)
 =\prod_{\sigma\in\Irr(Q)}P_\sigma(z)^{\dim\sigma}.
\end{equation}
Removing the simple scalar Perron factor and taking the normalized radial
logarithm at \(z=\lambda^{-1}\) gives
\begin{equation}\label{eq:quotient-cover-c-logarithm}
 \log C(\widetilde A_Q)=\log C(A)+
 \sum_{\substack{\sigma\in\Irr(Q)\\\sigma\ne\mathbf1}}
   (\dim\sigma)L_{\chi_\sigma}(\lambda^{-1}).
\end{equation}
Replacing \(F\) by \(L\) in
\eqref{eq:normal-subgroup-fixed-coordinate} and using
\eqref{eq:quotient-cover-c-logarithm} proves
\eqref{eq:quotient-artin-candidate}.

Although \(L_{\mathbf1_H}(\lambda^{-1})\) and
\(F_{\mathbf1_H}(\lambda^{-1})\) diverge separately, their radial difference
does not.  For \(0\le z<\lambda^{-1}\), the primitive-orbit expansions give
\begin{align}
 L_{\mathbf1_H}(z)-F_{\mathbf1_H}(z)
 &=\sum_{\omega\in\cP}\sum_{r\ge1}
   \frac{\mathbf1_H(g_\omega^r)-\mathbf1_H(g_\omega)}{r}
   z^{r|\omega|} \notag\\
 &=\sum_{\substack{\omega\in\cP\\q_H(\omega)>1}}
   \sum_{k\ge1}\frac{z^{kq_H(\omega)|\omega|}}
                       {kq_H(\omega)}.
 \label{eq:quotient-radial-difference}
\end{align}
Indeed, an orbit in \(H\) contributes zero, while for an orbit outside
\(H\), \(g_\omega^r\in H\) holds precisely when \(q_H(\omega)\mid r\).
Since \(q_H(\omega)\ge2\) in the last sum and the number of primitive
length-\(n\) orbits is \(O(\lambda^n/n)\), the endpoint series is absolutely
convergent.  Equations
\eqref{eq:normal-subgroup-fixed-coordinate}--
\eqref{eq:quotient-cover-c-logarithm} identify its endpoint with
\(\log(K_H/K_H^{\mathrm{Art}})\).  Summing the inner geometric logarithm
proves \eqref{eq:quotient-artin-positive-error}, hence
\eqref{eq:quotient-cover-splitting}.  Every summand in
\eqref{eq:quotient-artin-positive-error} is non-negative and is positive
exactly for an orbit outside \(H\), proving the last assertion.
\end{proof}

\begin{corollary}[Universal failure for the identity class]
\label{cor:universal-identity-artin-failure}
If \(G\ne\{e\}\) and the strict twisted gap holds, then
\begin{equation}\label{eq:identity-class-universal-error}
 \frac{K_{\{e\}}}{K_{\{e\}}^{\mathrm{Art}}}
 =\prod_{\substack{\omega\in\cP\\g_\omega\ne e}}
   \left(1-\lambda^{-\operatorname{ord}(g_\omega)|\omega|}\right)^{
      -1/\operatorname{ord}(g_\omega)}>1.
\end{equation}
Thus the Artin-coordinate formula strictly underestimates the identity-class
constant in every non-trivial strict-gap extension.
\end{corollary}

\begin{proof}
Take \(H=\{e\}\) in \Cref{thm:quotient-cover-splitting}.  It remains only to
exclude the possibility that every primitive orbit has identity holonomy.
If that occurred, then every period-\(n\) holonomy would be the identity, so
\(\Tr(A_\rho^n)=(\dim\rho)\Tr(A^n)\) for every irreducible \(\rho\) and every
\(n\ge1\).  Equality of all power sums would make the spectrum of \(A_\rho\)
equal to \(\dim\rho\) copies of the spectrum of \(A\).  A non-trivial finite
group has a non-trivial irreducible representation, for which this would give
\(\rad(A_\rho)=\lambda\), contrary to the strict gap.
\end{proof}

The determinant factorization
\eqref{eq:quotient-regular-determinant-factorization} is standard in
graph-covering Artin theory \cite{StarkTerras1996FiniteGraphsCoverings}.  The
need for auxiliary local correction factors has an arithmetic analogue in
Mertens theorems for Chebotarev prime sets
\cite{ArangoPinerosKeliherKeyes2022ChebotarevMertens}.  The assertion here is
the dynamical critical fixed-label factor
\eqref{eq:quotient-cover-splitting}, expressed geometrically through primitive
base orbits whose quotient lifts close only after more than one traversal.

\subsection{Harmonic renormalization and the complete algebraic jet}

Let \(\mathscr U\subseteq G\) be conjugacy invariant, and define
\begin{equation}\label{eq:union-class-product}
 \alpha_{\mathscr U}:=\frac{|\mathscr U|}{|G|},\qquad
 P_{\mathscr U}(N):=
 \prod_{\substack{\omega\in\cP\\|\omega|\le N,\ g_\omega\in\mathscr U}}
   (1-\lambda^{-|\omega|}),\qquad
 K_{\mathscr U}:=\prod_{C\subset\mathscr U}K_C.
\end{equation}
Here and below \(C\subset\mathscr U\) ranges over conjugacy classes.  Put
\begin{equation}\label{eq:refined-exponential-rate}
 \theta:=\max_{\substack{\nu\in\spec(A)\\\nu\ne\lambda}}|\nu|,
 \qquad
 q_0:=\max\left\{\frac{\theta}{\lambda},
                   \frac{\eta}{\lambda},\lambda^{-1/2}\right\}<1,
\end{equation}
with the maximum over an empty spectrum equal to zero.

\begin{theorem}[Exponentially accurate harmonic renormalization]
\label{thm:harmonic-renormalization}
Under the hypotheses of \Cref{thm:correct-product}, for every conjugacy
invariant \(\mathscr U\subseteq G\) and every \(q_0<\xi<1\),
\begin{equation}\label{eq:harmonic-renormalized-product}
 P_{\mathscr U}(N)=K_{\mathscr U}N^{-\alpha_{\mathscr U}}
 \exp\left\{-\alpha_{\mathscr U}
   (H_N-\log N-\gamma)\right\}
 \left(1+O_\xi(\xi^N)\right).
\end{equation}
Consequently,
\begin{equation}\label{eq:harmonic-logarithmic-jet}
 \log\frac{P_{\mathscr U}(N)N^{\alpha_{\mathscr U}}}{K_{\mathscr U}}
 \sim-\frac{\alpha_{\mathscr U}}{2N}
 +\alpha_{\mathscr U}\sum_{j\ge1}
    \frac{B_{2j}}{2jN^{2j}},
\end{equation}
and, in particular,
\begin{align}
 P_{\mathscr U}(N)
 =K_{\mathscr U}N^{-\alpha_{\mathscr U}}
 \bigg(&1-\frac{\alpha_{\mathscr U}}{2N}
 +\frac{\alpha_{\mathscr U}(3\alpha_{\mathscr U}+2)}{24N^2}
 \notag\\
 &-\frac{\alpha_{\mathscr U}^2(\alpha_{\mathscr U}+2)}{48N^3}
 +O(N^{-4})\bigg).
 \label{eq:harmonic-product-jet-three}
\end{align}
Thus all algebraic-order coefficients depend on the system only through
\(\alpha_{\mathscr U}\); the remaining system dependence is in
\(K_{\mathscr U}\) and an exponentially small error.
\end{theorem}

\begin{proof}
Let \(p_{n,\mathscr U}\) count the primitive length-\(n\) orbits with label
in \(\mathscr U\), and set
\begin{equation}\label{eq:harmonic-error-sequences}
 a_n:=p_{n,\mathscr U}\lambda^{-n}
      -\frac{\alpha_{\mathscr U}}{n},
 \qquad
 b_n:=p_{n,\mathscr U}\sum_{r\ge2}\frac{\lambda^{-rn}}{r}.
\end{equation}
The scalar M\"obius formula in the proof of \Cref{lem:class-mertens}, the
class-indicator expansion, and
\eqref{eq:explicit-primitive-trace-bound} imply, for every \(q_0<\xi<1\),
\begin{equation}\label{eq:harmonic-error-bounds}
 a_n=O_\xi(\xi^n),\qquad b_n=O(\lambda^{-n}/n).
\end{equation}
The divisor factor in the M\"obius bounds is absorbed by the strict
inequality \(q_0<\xi\).

Taking real logarithms in \eqref{eq:union-class-product} gives
\[
 \log P_{\mathscr U}(N)
 =-\alpha_{\mathscr U}H_N-
   \sum_{n\le N}(a_n+b_n).
\]
On letting \(N\to\infty\) and comparing with the constant supplied by
\Cref{thm:correct-product}, one obtains
\(\log K_{\mathscr U}=-\alpha_{\mathscr U}\gamma-
\sum_{n\ge1}(a_n+b_n)\).  Hence the exact identity
\begin{equation}\label{eq:exact-harmonic-remainder}
 \log\frac{P_{\mathscr U}(N)N^{\alpha_{\mathscr U}}}{K_{\mathscr U}}
 =-\alpha_{\mathscr U}(H_N-\log N-\gamma)
  +\sum_{n>N}(a_n+b_n).
\end{equation}
The last sum is \(O_\xi(\xi^N)\) by
\eqref{eq:harmonic-error-bounds}; exponentiation proves
\eqref{eq:harmonic-renormalized-product}.  Finally, Euler--Maclaurin gives
\[
 H_N-\log N-\gamma
 \sim \frac{1}{2N}-\sum_{j\ge1}\frac{B_{2j}}{2jN^{2j}}.
\]
Substitution proves \eqref{eq:harmonic-logarithmic-jet}; exponentiating its
first terms gives \eqref{eq:harmonic-product-jet-three}.
\end{proof}

For Ihara prime cycles, Hasegawa and Saito obtained the scalar all-order
Bernoulli expansion for the second graph-theoretic Mertens sum
\cite{HasegawaSaito2016GraphMertens}.  In the aperiodic case their scalar
coefficients contain the \(\alpha_{\mathscr U}=1\) harmonic jet above, so no
priority is claimed here for those scalar Bernoulli coefficients.  The added
content of \Cref{thm:harmonic-renormalization} is the exact harmonic
renormalization, with exponential remainder, simultaneously for every
fixed-label union of Frobenius classes under the strict twisted gap.
